% TODO make hw stylesheet
\documentclass[11pt]{scrartcl}
\usepackage[a4paper, total={6in, 8in}]{geometry}
\usepackage[usenames,dvipsnames,svgnames]{xcolor}
\usepackage[shortlabels]{enumitem}
\usepackage[framemethod=TikZ]{mdframed}
\usepackage{amsmath,amssymb,amsthm}
\usepackage[colorlinks]{hyperref}
\usepackage{multicol}
\usepackage{tabularx}
\usepackage{stmaryrd}

\hypersetup{
	colorlinks=true,
	linkcolor=blue,
	filecolor=magenta,
	urlcolor=magenta,
	pdftitle={MAT 528 HW}
}

\usepackage{microtype}
\usepackage{mathtools}
\usepackage[headsepline]{scrlayer-scrpage}
\usepackage{thmtools}
\usepackage[textsize=footnotesize]{todonotes}

\mdfdefinestyle{mdbluebox}{roundcorner=10pt,innerbottommargin=9pt,
	linecolor=blue,backgroundcolor=TealBlue!5,}
\declaretheoremstyle[headfont=\sffamily\bfseries\color{MidnightBlue},
mdframed={style=mdbluebox},]{thmbluebox}

\mdfdefinestyle{mdbloobox}{frametitlefont=\bfseries,innerbottommargin=8pt,
	nobreak=true,backgroundcolor=TealBlue!5,linecolor=blue,}
\declaretheoremstyle[headfont=\bfseries\color{MidnightBlue},
mdframed={style=mdbloobox},headpunct={\\[3pt]},postheadspace=0pt,]{thmbloobox}

\mdfdefinestyle{mdredbox}{frametitlefont=\bfseries,innerbottommargin=8pt,
	nobreak=true,backgroundcolor=Salmon!5,linecolor=RawSienna,}
\declaretheoremstyle[headfont=\bfseries\color{RawSienna},
mdframed={style=mdredbox},headpunct={\\[3pt]},postheadspace=0pt,]{thmredbox}

\mdfdefinestyle{mdgreenbox}{linecolor=ForestGreen,backgroundcolor=ForestGreen!5,
	linewidth=2pt,rightline=false,leftline=true,topline=false,bottomline=false,}
\declaretheoremstyle[headfont=\bfseries\sffamily\color{ForestGreen!70!black},
mdframed={style=mdgreenbox},headpunct={ --- },]{thmgreenbox}

\mdfdefinestyle{mdorangebox}{linecolor=RedOrange,backgroundcolor=RedOrange!5,
	linewidth=2pt,rightline=false,leftline=true,topline=false,bottomline=false,}
\declaretheoremstyle[headfont=\bfseries\sffamily\color{RedOrange!70!black},
mdframed={style=mdorangebox},headpunct={ --- },]{thmorangebox}

\mdfdefinestyle{mdblackbox}{linecolor=black,backgroundcolor=RedViolet!5!gray!5,
	linewidth=3pt,rightline=false,leftline=true,topline=false,bottomline=false,}
\declaretheoremstyle[mdframed={style=mdblackbox}]{thmblackbox}

\declaretheoremstyle[spaceabove=6pt,spacebelow=6pt]{basehead}

\declaretheorem[style=basehead,name=Problem]{problem}
\declaretheorem[style=basehead,name=Problem]{problem*}
\declaretheorem[style=thmbloobox,name=Setup]{setup} % for config geo
\declaretheorem[style=thmredbox,name=Example,sibling=problem]{example}
\declaretheorem[style=thmredbox,name=$\bigstar$ Problem,sibling=problem]{niceprob}
\declaretheorem[style=thmbluebox,name=Theorem,sibling=problem]{theorem}
\declaretheorem[style=thmbluebox,name=Corollary,sibling=theorem]{corollary}
\declaretheorem[style=thmbluebox,name=Proposition,sibling=theorem]{prop}
\declaretheorem[style=thmbluebox,name=Lemma,numberwithin=problem]{lemma}
\declaretheorem[style=thmbluebox,name=Theorem,numbered=no]{theorem*}
\declaretheorem[style=thmbluebox,name=Lemma,numbered=no]{lemma*}
\declaretheorem[style=thmgreenbox,name=Claim,numberwithin=problem]{claim}
\declaretheorem[style=thmgreenbox,name=Claim,numbered=no]{claim*}
\declaretheorem[style=thmblackbox,name=Remark,numberwithin=problem]{remark}
\declaretheorem[style=thmblackbox,name=Remark,numbered=no]{remark*}
\declaretheorem[style=thmgreenbox,name=Definition,sibling=theorem]{definition}
\declaretheorem[style=thmgreenbox,name=Definition,numbered=no]{definition*}
\declaretheorem[style=thmorangebox,name=Warning,sibling=theorem]{warning}
\declaretheorem[style=thmorangebox,name=Warning,numbered=no]{warning*}
\declaretheorem[style=thmorangebox,name=Note,numbered=no]{note}
\declaretheorem[style=thmblackbox,name=Exercise,sibling=problem]{exercise}
\declaretheorem[style=thmblackbox,name=Exercise,numbered=no]{exercise*}
\declaretheorem[style=thmblackbox,name=Question,sibling=problem]{question}
\declaretheorem[style=thmblackbox,name=Question,numbered=no]{question*}

\addtolength{\textheight}{3.14cm}
\providecommand{\ol}{\overline}
\providecommand{\ceil}[1]{\left\lceil #1 \right\rceil}
\providecommand{\floor}[1]{\left\lfloor #1 \right\rfloor}
\newcommand{\dd}{\mathrm{d}}
\providecommand{\ul}{\underline}
\providecommand{\eps}{\varepsilon}
\providecommand{\half}{\frac{1}{2}}
\providecommand{\dang}{\measuredangle} %% Directed angle
\providecommand{\CC}{\mathbb C}
\providecommand{\FF}{\mathbb F}
\providecommand{\NN}{\mathbb N}
\providecommand{\QQ}{\mathbb Q}
\providecommand{\RR}{\mathbb R}
\providecommand{\ZZ}{\mathbb Z}
\providecommand{\dg}{^\circ}
\providecommand{\ii}{\item}
\providecommand{\setdiff}{\smallsetminus}
\providecommand{\alert}[1]{{\sffamily\textbf{\textcolor{blue}{#1}}}}
\providecommand{\inv}{^{-1}}
\providecommand{\mb}[1]{\mathbf{#1}}

\reversemarginpar

\newenvironment{soln}{\begin{proof}[Solution]}{\end{proof}}
\newenvironment{walkthrough}{\noindent \textbf{Walkthrough.}}{\medskip}
\newlist{walk}{enumerate}{3}
\setlist[walk]{label=\bfseries (\alph*)}

\providecommand{\pow}{\operatorname{Pow}}
\providecommand{\vocab}[1]{{\textbf{\textcolor{ForestGreen}{#1}}}}
\providecommand{\lcm}{\operatorname{lcm}}
\providecommand{\abs}[1]{\left|#1\right|}

\usepackage{asymptote}
\begin{asydef}
	defaultpen(fontsize(10pt));
	unitsize(1cm);
	size(8cm); // set a reasonable default
	usepackage("amsmath");
	usepackage("amssymb");
	settings.tex="pdflatex";
	settings.outformat="pdf";
	// Replacement for olympiad+cse5 which is not standard
	import geometry;
	// recalibrate fill and filldraw for conics
	void filldraw(picture pic = currentpicture, conic g, pen fillpen=defaultpen, pen drawpen=defaultpen)
	{ filldraw(pic, (path) g, fillpen, drawpen); }
	void fill(picture pic = currentpicture, conic g, pen p=defaultpen)
	{ filldraw(pic, (path) g, p); }
	// some geometry
	pair foot(pair P, pair A, pair B) { return foot(triangle(A,B,P).VC); }
	pair orthocenter(pair A, pair B, pair C) { return orthocentercenter(A,B,C); }
	pair centroid(pair A, pair B, pair C) { return (A+B+C)/3; }
	// cse5 abbreviations
	path CP(pair P, pair A) { return circle(P, abs(A-P)); }
	path CR(pair P, real r) { return circle(P, r); }
	pair IP(path p, path q) { return intersectionpoints(p,q)[0]; }
	pair OP(path p, path q) { return intersectionpoints(p,q)[1]; }
	path Line(pair A, pair B, real a=0.6, real b=a) { return (a*(A-B)+A)--(b*(B-A)+B); }
	// cse5 more useful functions
	picture CC() {
		picture p=rotate(0)*currentpicture;
		currentpicture.erase();
		return p;
	}
	pair MP(Label s, pair A, pair B = plain.S, pen p = defaultpen) {
		Label L = s;
		L.s = "$"+s.s+"$";
		label(L, A, B, p);
		return A;
	}
	pair Drawing(Label s = "", pair A, pair B = plain.S, pen p = defaultpen) {
		dot(MP(s, A, B, p), p);
		return A;
	}
	path Drawing(path g, pen p = defaultpen, arrowbar ar = None) {
		draw(g, p, ar);
		return g;
	}
\end{asydef}

\usepackage{mathrsfs}

\ihead{\footnotesize MAT 528 HW01}
\ohead{\footnotesize September 4, 2025}

\title{\vspace{-2em}MAT 528 HW01}
\author{\large Aditya Pahuja}
\date{\large Due 09/04}
\begin{document}
\maketitle

\begin{problem*}
    [Munkres 13.4]
    \begin{enumerate}[(a)]
        \ii If $\{\mathcal T_\alpha\}$ is a family of topologies on $X$,
	show that $\bigcap \mathcal T_\alpha$ is a topology on $X$.
	Is $\bigcup\mathcal T_\alpha$ is a topology on $X$?
	\ii Let $\{\mathcal T_\alpha\}$ be a family of topologies on $X$.
	Show that there is a unique smallest topology on $X$ containing
	all the collections $\{\mathcal T_\alpha\}$,
	and a unique largest topology contained in all the $\mathcal T_\alpha$.
	\ii If $X = \{a,b,c\}$, let
	\begin{equation*}
	    \mathcal T_1 = \{\varnothing, X, \{a\}, \{a,b\}\}\quad\text{and}\quad
	    \mathcal T_2 = \{\varnothing, X, \{a\}, \{b,c\}\}.
	\end{equation*}
	Find the smallest topology containing $\mathcal T_1$ and $\mathcal T_2$,
	and the largest topology contained in $\mathcal T_1$ and $\mathcal T_2$.
    \end{enumerate}
\end{problem*}

\begin{soln}
    (a) We check that $\bigcap_\alpha\mathcal T_\alpha$ satisfies
    the three axioms for a collection of subsets on $X$ to be a topology.
    \begin{itemize}
        \ii Each of the $\mathcal T_\alpha$ contains $X$ and the empty set,
	so the intersection of the $\mathcal T_\alpha$
	also contains $X$ and $\varnothing$.

        \ii Suppose that $\{U_\beta\}_{\beta\in B}$ is a collection
        of sets in $\bigcap_\alpha\mathcal T_\alpha$.
        Then $\{U_\beta\}_{\beta\in B}\subseteq \mathcal T_\alpha$
	for each $\alpha\in A$,
        so $\bigcup_{\beta\in B}U_\beta$ is in each $\mathcal T_\alpha$,
        hence $\bigcup_{\beta\in B}U_\beta$ is in the intersection of the
        $\mathcal T_\alpha$.

	\ii Finally, suppose that $U_1$, \dots, $U_n$ are $n$ sets in the intersection
	of the $\mathcal T_\alpha$. Then these $n$ sets are contained
	in $\mathcal T_\alpha$ for all $\alpha\in A$, so
	their intersection is in each $\mathcal T_\alpha$,
	hence $U_1\cap\cdots\cap U_n$ lies in $\bigcap_{\alpha}\mathcal T_\alpha$.
    \end{itemize}
    Therefore the intersection of the $\mathcal T_\alpha$ is a topology on $X$.

    On the other hand, unions of topologies need not be topologies.
    For example, $\mathcal T_1 = \{\varnothing, X, \{a\}\}$
    and $\mathcal T_2 = \{\varnothing, X, \{b\}\}$ are topologies
    on $X = \{a,b,c\}$, but their union $\mathcal T_1\cup\mathcal T_2$
    is not a topology because it contains $\{a\}$ and $\{b\}$
    but not their union $\{a, b\}$.\\

    (b) Let $\mathcal T_L$ be the intersection of all topologies on $X$
    containing $\bigcup_\alpha\mathcal T_\alpha$.
    Clearly this is a topology containing all the $\mathcal T_\alpha$
    that is minimal with respect to inclusion,
    since if $\mathcal T'_L$ is a topology containing all the $\mathcal T_\alpha$
    then it is part of the intersection that generated $\mathcal T_L$,
    hence $\mathcal T_L\subseteq\mathcal T'_L$.
    This topology is also the \emph{unique} minimal topology
    containing all the $\mathcal T_\alpha$ because
    if $\mathcal T''_L$ is another minimal topology
    then by minimality $\mathcal T_L$ and $\mathcal T''_L$ contain each other.

    Similarly, the largest topology contained in all the $\mathcal T_\alpha$
    is $\bigcap_\alpha \mathcal T_\alpha$. This is a topology by part (a)
    and it is easy to see that it is the unique largest one
    because $\bigcap_\alpha\mathcal T_\alpha$ is straight up
    the largest set contained in all of the $\mathcal T_\alpha$.\\

    (c) The smallest topology containing both is
    \begin{equation*}
	\{\varnothing, X, \{a\}, \{b\},
	\{a,b\}, \{b,c\}\}
    \end{equation*}
    since it must contain $\mathcal T_1\cup\mathcal T_2$ of course,
    and in addition $\{b\} = \{a,b\}\cap\{a,c\}$ is in this topology as well,
    which accounts for all $6$ sets.
    The largest topology contained in $\mathcal T_1$ and $\mathcal T_2$
    is just their intersection
    \begin{equation*}
	\{\varnothing, X, \{a\}\}.\qedhere
    \end{equation*}
\end{soln}

\begin{problem*}
    [Munkres 13.5]
    Show that if $\mathcal A$ is a basis for a topology on $X$,
    then the topology generated by $\mathcal A$ equals
    the intersections of all topologies on $X$ that contain $\mathcal A$.
    Prove the same if $\mathcal A$ is a subbasis.
\end{problem*}

\begin{soln}
    By the previous problem, the topology generated by $\mathcal A$
    must contain the intersection of all topologies containing $\mathcal A$
    because the intersection is minimal with respect to inclusion
    among topologies containing $\mathcal A$.
    Conversely, any set in the topology $\mathcal T$ generated by $\mathcal A$
    is the union of some sets in $\mathcal A$. Since topologies
    are closed under arbitrary unions, every topology containing $\mathcal A$
    contains arbitrary unions of elements of $\mathcal A$,
    so all sets in $\mathcal T$ are contained in
    every topology containing $\mathcal A$,
    which means the topology generated by $\mathcal A$ is contained by
    the intersection of all topologies containing $\mathcal A$.

    If $\mathcal A$ is a subbasis then again the topology generated by $\mathcal A$
    contains the intersection of all topologies containing $\mathcal A$.
    For the converse, we note that open sets in the topology
    generated by $\mathcal A$ are arbitrary unions of finite intersections
    of sets in $\mathcal A$. Topologies are closed under finite intersections
    and arbitrary unions, so any topology containing $\mathcal A$
    contains all finite intersections in $\mathcal A$
    and thus all arbitrary unions of finite intersections of sets in $\mathcal A$.
    That is, every topology on $X$ containing $\mathcal A$
    contains the topology generated by $\mathcal A$,
    i.e. the intersection of all topologies on $X$ containing $\mathcal A$
    contains the topology generated by $\mathcal A$.
\end{soln}

\begin{problem*}
    [Munkres 13.8a]
    Apply Lemma 13.2 to show that the countable collection
    \begin{equation*}
	\mathcal B = \{(a,b) : a<b,\; \text{$a$ and $b$ are rational}\}
    \end{equation*}
    is a basis that generates the standard topology on $\RR$.
\end{problem*}

\begin{soln}
    Let $x$ be an element of $\RR$ and $U\subseteq\RR$
    an open set containing $x$. Then, $U$ contains some interval
    $(x-\eps,x+\eps)$ centered at $x$, since the standard topology on $\RR$
    is defined with open intervals as a basis.
    Pick rational numbers $a\in (x-\eps,x)$ and $b\in(x,x+\eps)$;
    then $x\in(a,b)\subseteq U$ and $(a,b)\in\mathcal B$.
    By Lemma 13.2 this is enough to imply that $\mathcal B$
    is a basis for the standard topology on $x$.
\end{soln}

\begin{problem*}
    [Munkres 16.3]
    Consider the set $Y = [-1,1]$ as a subspace of $\RR$.
    Which of the following sets are open in $Y$?
    Which are open in $\RR$?
    \begin{align*}
	A &= \{x : \tfrac12 < |x| < 1\} \\
	B &= \{x : \tfrac12 < |x| \le 1\} \\
	C &= \{x : \tfrac12 \le |x| < 1\} \\
	D &= \{x : \tfrac12 \le |x| \le 1\} \\
	E &= \{x : 0 < |x| < 1\text{ and }1/x\notin\ZZ_+\}
    \end{align*}
\end{problem*}

\begin{soln}
    \begin{itemize}
        \ii $A$ is open in $\RR$ because it is the union of intervals
	$(-1,-\half)\cup(\half,1)$, and it is open in $Y$
	because $A = A\cap Y$ where $A$ is open in $\RR$.
	\ii $B$ is not open in $\RR$ because any interval centered at $1$
	must contain some points greater than $1$,
	which means said interval is not inside $B$.
	On the other hand, $B$ is open in $Y$ because
	$B = ((-2,\half)\cup(\half,2))\cap Y$.
	\ii $C$ is not open in $\RR$ because any interval centered at $\half$
	contains points outside $C$ (specifically some points just below $\half$).
	It is also not open in $Y$ because of the same reason,
	since if $C = U\cap Y$ for some $U\subseteq\RR$ then
	$U$ would have the same behavior around $\half$,
	so $U$ isn't open in $\RR$.
	\ii $D$ is not open in $\RR$ again because any interval around $\half$
	must contain some points greater than $1$
	so the interval isn't contained within $D$,
	It is also again not open in $Y$ because
	if $D = U\cap Y$ for some $U\subseteq\RR$ then $U$ cannot be open in $\RR$
	because it has the same behavior around $\half$.
	\ii $E$ is open in $\RR$ because it is the union of open intervals
	\begin{equation*}
	    \bigcup_{n=1}^\infty \left(\frac1{n+1},\frac1n\right)\cup
	    \bigcup_{n=1}^\infty \left(-\frac1{n},-\frac1{n+1}\right)
	\end{equation*}
	and it is open in $Y$ because each of the intervals in this union
	is strictly contained within $Y$, hence also open in $Y$.
    \end{itemize}
\end{soln}

\begin{problem*}
    [Munkres 16.4]
    A map $f\colon X\to Y$ is said to be an \vocab{open map} if
    for every open set $U$ of $X$, the set $f(U)$ is open in $Y$.
    Show that $\pi_1\colon X\times Y\to X$ and $\pi_2\colon X\times Y\to Y$
    are open maps.
\end{problem*}

\begin{soln}
    Recall that a basis for the product topology on $X$
    is the collection of sets $\{U_X\times U_Y : U_X\subseteq X,\, U_Y\subseteq Y\}$.
    If $U\subseteq X\times Y$ is open, then for each $p\in U$
    there is a basis element $U_{X,p}\times U_{Y,p}$ containing $p$
    with $U_{X,p}\subseteq X$ and $U_{Y,p}\subseteq Y$.
    When projected onto $X$, we get
    \begin{equation*}
	\pi_1(U) = \pi_1\left(\bigcup_{p\in U}U_{X,p}\times U_{Y,p}\right)
	= \bigcup_{p\in U}\pi_1(U_{X,p}\times U_{Y,p})
	= \bigcup_{p\in U}U_{X,p}\subseteq X
    \end{equation*}
    which is the union of open subsets of $X$, hence is open in $X$.
    Similarly
    \begin{equation*}
	\pi_2(U) = \bigcup_{p\in U}U_{Y,p}\subseteq Y
    \end{equation*}
    is open in $Y$ because it is the union of open subsets of $Y$.
\end{soln}

\begin{problem*}
    [Munkres 16.6]
    Show that the countable collection
    \begin{equation*}
	\{(a,b)\times(c,d) : a<b\text{ and }c<d\text{ and $a,b,c,d$ are rational}\}
    \end{equation*}
    is a basis for $\RR^2$.
\end{problem*}

\begin{soln}
    This is an immediate corollary of the fact that
    if $\mathcal B_X$ and $\mathcal B_Y$ are bases for
    topologies on $X$ and $Y$, then
    \begin{equation*}
	\mathcal B_{X\times Y} =
	\{ U\times V : U\in\mathcal B_X,\,V\in\mathcal B_Y\}
    \end{equation*}
    is a basis for the product topology on $X\times Y$ ---
    in this case, $X = Y = \RR^2$
    and $\mathcal B_X = \mathcal B_Y$ is the basis consisting of
    intervals with rational endpoints, as per the result of
    Problem 3 (Munkres 13.8a).
\end{soln}

\pagebreak
\begin{problem*}
    [Munkres 16.8]
    If $L$ is a straight line in the plane, describe the topology $L$ inherits
    as a subspace of $\RR_\ell\times\RR$ and as a subspace of $\RR_\ell\times\RR_\ell$.
\end{problem*}

\begin{soln}
    The basis elements of $\RR_\ell\times\RR$ are of the form
    $[a,b)\times(c,d)$ for reals $a<b$, $c<d$;
    intersecting them with $L$ gives a basis for the subspace topology on $L$.
    If $L$ is vertical, then $L$ is homeomorphic to $\RR$ with the standard topology
    because the intersection of a basis set with $L$ is either empty
    or a set of the form $\{x_0\}\times(c,d)$ where $x_0$ is the $x$-coordinate
    of the points on the line.
    If $L$ isn't vertical, then $L$ is homeomorphic to $\RR_\ell$
    because the segment along which $L$ intersects a basis set
    either has its left endpoint or neither endpoint,
    and the latter case covered by the former because $\RR_\ell$ is finer
    than $\RR$ with the standard topology.

    The basis elements of $\RR_\ell\times\RR_\ell$ are of the form
    $[a,b)\times[c,d)$ for real $a<b$, $c<d$.
    If $L$ has negative slope then for any point $(x,y)\in L$,
    \begin{equation*}
	L\cap[x,x+1)\times[y,y+1) = (x,y)
    \end{equation*}
    is open in $L$, so $L$ just gets the discrete topology.
    Otherwise, $L$ is homeomorphic to $\RR_\ell$
    because every intersection of $L$ with a basis set
    will be a segment with its left endpoint but not its right endpoint.
\end{soln}

\end{document}
